Fix a feasible \((p,w)\) and a schedule in Definition~\(\mathrm{def{:}ellipsoidal\text{-}coefficient\text{-}class}\).  The low-order expansion and Assumption~\(\mathrm{ass{:}common\text{-}uniform\text{-}rollout}\) give
\[
\mathbb E_U M(r)=\bar c_{\varnothing}
+\sum_{S\in\mathcal S_{n,\beta}}a_Sr^{|S|}.
\tag{37}
\]
By Definition~\(\mathrm{def{:}unbiased\text{-}weight\text{-}class}\), \(\sum_tw_t=0\) and \(\sum_tw_tp_t^d=1\) for \(1\le d\le\beta\).  Applying these equations term by term in (37) yields
\[
\mathbb E_U\widehat\tau_{p,w}=\sum_{S\in\mathcal S_{n,\beta}}a_S.
\tag{38}
\]
Assumption~\(\mathrm{ass{:}low\text{-}order\text{-}schedule}\) separately gives the causal contrast
\[
\tau=\frac1n\sum_{i\in I_n}\{Y_i(\mathbf1)-Y_i(\mathbf0)\}
=\sum_{S\in\mathcal S_{n,\beta}}a_S.
\tag{39}
\]
Thus (38)--(39) prove design unbiasedness rather than defining it.

Lemma~\(\mathrm{lem{:}common\text{-}uniform\text{-}monomial\text{-}covariance}\) gives
\[
\operatorname{Var}_U(\widehat\tau_{p,w})=a^{\mathsf T}G(p,w)a.
\]
Since every \(\omega_d>0\), set \(b=\Omega^{-1/2}a\).  The ellipsoid assumption gives \(\|b\|_2^2\le B^2/n\), and the Rayleigh inequality gives
\[
\begin{aligned}
a^{\mathsf T}G(p,w)a
&=b^{\mathsf T}\Omega^{1/2}G(p,w)\Omega^{1/2}b\\
&\le\frac{B^2}{n}\lambda_{\max}\{\Omega^{1/2}G(p,w)\Omega^{1/2}\}\\
&=V_{\mathrm{rob}}(p,w),
\end{aligned}
\tag{40}
\]
where the last equality uses Lemma~\(\mathrm{lem{:}subset\text{-}lattice\text{-}block\text{-}reduction}\) and Definition~\(\mathrm{def{:}variance\text{-}certificate}\).

Suppose \(V_{\mathrm{rob}}(p,w)>0\), put \(X=\widehat\tau_{p,w}-\tau\), and set \(R=\sqrt{V_{\mathrm{rob}}(p,w)/\alpha}\).  Equations (38)--(40) give \(\mathbb E_UX=0\) and \(\mathbb E_UX^2\le V_{\mathrm{rob}}(p,w)\).  The pointwise inequality
\[
X^2\ge R^2\mathbf1\{|X|>R\}
\]
implies
\[
\Pr_U(|X|>R)\le\frac{\mathbb E_UX^2}{R^2}\le\alpha.
\]
Taking complements and using Definition~\(\mathrm{def{:}chebyshev\text{-}interval}\) proves the coverage claim.  If \(V_{\mathrm{rob}}(p,w)=0\), (40) gives \(\mathbb E_UX^2=0\); nonnegativity of \(X^2\) implies \(X=0\) almost surely.  The singleton interval therefore covers with probability one.